\section{Parameter absorption and objective}
The explicit selection rule is
\(B_0=\lceil2(512+1)\rceil=1026\),
\(D_0=\lceil2(256+1)\rceil=514\).
For \(x\ge32768\), \(x\ge2B_0\), \(x^2\ge4(D_0+1)\), and
\[
 \frac12\le S\le1-\frac{514}{x^2},\qquad
 \frac{H^2}{2\ell}\le N,\qquad 1-\beta^2\ge\frac{1026}{x}.
\]
For the upper estimate on \(S\), put
\(q=(\ell-x^2)/x^2\le1/4\); then
\(1-(1+q)^{-2}\ge q\ge514/x^2\).
The square bracket in \eqref{structuredineq} is therefore at most \(-1/x\).
The elementary exponential bound gives
\[
 64e^{-x/24}\le\frac{64\,24^5}{x^4}\le x^{-2},
 \qquad 64\,24^5<2^{30}\le x^2.
\]
Thus all short and medium budgets are at most
\(1-257/x^2-w/x<1\), and Lemma~\ref{structured} covers the other lags.
The construction is feasible for every integer \(N\ge2^{60}\).

Positive part adds mass, while the cap changes only \(\Gamma\).
Using the BV quadrature bound and \(\int b=-1/6\) gives
\[
 \sum_j\psi_j\ge\delta H
 \left(U-\frac13-2\int_U^\infty|b|\right)-26\delta.
\]
The atoms have total mass \(\beta\delta H\) and \(HU=N-\ell\).
Consequently
\begin{equation}
 \kappa\ge\delta N-1+\frac23\delta H
 -\frac{1026\delta H}{x}-26\delta-12\delta H e^{-U/4}. \label{mass}
\end{equation}
Here \(\delta N\ge x^2-515\),
\(\delta H\ge\sqrt2x-1\),
\(1026\delta H/x<2052\), and \(26\delta<1\).
Also \(U\ge x/2\), and
\(e^{-U/4}\le24\,8^4/x^4\le x^{-2}\), so the last loss is at most
\(24/x<1\). The losses in \eqref{mass} total at most 2571.
The declared constant 2572 includes one further unit of slack.

\section{An attributed uniform upper companion}
Let \(\phi(t)=4/3-2|t|+(2/3)|t|^3\) on \(|t|\le1\), zero otherwise.
It is the autocorrelation of a normalized ramp. Thus
\(A_T(i,j)=\phi((i-j)/T)\) is PSD. For any feasible moment pair,
\[
 f^{\mathsf T}A_T f\le\operatorname{tr}(A_T\Gamma)
 \le\frac43\kappa+2\sum_{d\ge1}\phi(d/T)\le\frac43\kappa+T.
\]
The last estimate uses the decreasing nonnegative kernel and its unit
integral, and holds for real \(T>0\).
Chen's finite-positive-measure comparison \cite[Lemma 4.4]{chen}, applied to
\(\mu=\sum_j f_j\delta_{j/T}\) on \([0,N/T]\), yields
\begin{equation}
 \kappa^2\le\left(\frac NT+\frac23+200e^{-aN/T}\right)
 \left(T+\frac43\kappa\right),\qquad a=\log(4/3),\quad 0<T\le N.
 \label{upper}
\end{equation}
This uses energy Cauchy--Schwarz with a finite signed reference measure
whose potential is exactly one on the closed interval. Its energy is at most
\(N/T+2/3+200e^{-aN/T}\). It imposes no integrality on \(\kappa\), and
no division by \(\kappa\) is required.

For completeness, a simple large-onset root comparison suffices here.
Set \(T=\sqrt2x^3\), \(g=2\sqrt2/3\), and
\(\eta=200e^{-ax/\sqrt2}\). Equation~\eqref{upper} says
\[
 P_\eta(\kappa)\le0,\quad
 P_\eta(z)=z^2-(gx+8/9+4\eta/3)z-x^4-gx^3-\sqrt2\eta x^3.
\]
At \(y=x^2+gx+1\), direct expansion gives
\(P_0(y)=(10/9)x^2+(g/9)x+1/9\).
Since \(a>1/4\) and \(\sqrt2<3/2\),
\[
 \eta\le200e^{-x/6}\le\frac{200\cdot24\cdot6^4}{x^4}\le x^{-2}
 \quad(x\ge32768).
\]
The error at \(y\) is at most \(4+2x<(10/9)x^2\). Hence
\(P_\eta(y)>0\). The quadratic has negative constant term and one positive
root, so \(\kappa<y\), proving the upper half of Theorem~\ref{main}.
Chen's sharper rational absorption gives this upper companion already for
\(N\ge4{,}600{,}000\); its all-real-mass replay is retained in the evidence
packet. The lower onset \(2^{60}\) dominates it.
